% =============================================================================
% CONCEPTUAL / TYPES-AND-VALUES POOL
%
% Four items, each targeting a specific fundamental misunderstanding rather
% than a language-trivia fact. All SML behavior verified in SML/NJ.
%
%   1. A function body does not run until the function is applied.
%   2. Static scope: a closure captures bindings at its definition site.
%   3. Equivalence is not runtime cost.
%   4. Mutation breaks referential transparency.
%
% Each is written as a self-contained \part so it can be dropped into either
% the Types-and-Values or Conceptual Questions section of either exam.
% =============================================================================


% -----------------------------------------------------------------------------
% 1. A function body does not run until applied.
%    Misunderstanding tested: students believe binding a `fn` (or defining a
%    fun) evaluates its body, so effects/exceptions inside fire "at definition".
% -----------------------------------------------------------------------------
\part[3]\hfill

  Consider the following sequence of declarations, evaluated in order:

  \begin{codeblock}
    val a = 1
    val f = fn () => 1 div 0
    val g = fn x => raise Fail "boom"
    val res = a + a
  \end{codeblock}

  Does evaluating these declarations raise an exception? Answer yes or no, and
  justify your answer in one sentence. If it does not raise, state the value of
  \code{res}.

  \begin{solutionorbox}[9em]
    No. Binding \code{f} and \code{g} does \textbf{not} run their bodies --- a
    function's body is only evaluated when the function is \textit{applied},
    and neither \code{f} nor \code{g} is ever applied. So no division by zero
    and no \code{raise} occurs, and \code{res} $\eeq$ \code{2}.

    (Contrast \code{val c = 1 div 0}, which is not a function and \textit{would}
    raise immediately, because its right-hand side is evaluated at binding
    time.)
  \end{solutionorbox}


% -----------------------------------------------------------------------------
% 2. Static scope: a closure captures the binding in effect at its definition,
%    not at its call. Misunderstanding: students think name lookup happens at
%    call time against the "current" binding.
% -----------------------------------------------------------------------------
\part[3]\hfill

  Consider the following declarations, evaluated in order:

  \begin{codeblock}
    fun makeAdder () =
      let
        val n = 42
      in
        fn y => n + y
      end

    val add = makeAdder ()
    val n = 1000
    val res = add 1
  \end{codeblock}

  What is the value of \code{res}? Justify your answer in one sentence.

  \begin{solutionorbox}[9em]
    \code{res} $\eeq$ \code{43}.

    The function \code{fn y => n + y} captures the binding of \code{n} in scope
    \textbf{where it was defined} --- the \code{val n = 42} inside
    \code{makeAdder} --- not whatever \code{n} happens to be bound to at the
    point where \code{add} is later called. The later \code{val n = 1000} is a
    different, unrelated binding and has no effect on \code{add}.
  \end{solutionorbox}


% -----------------------------------------------------------------------------
% 3. Equivalence is not runtime cost.
%    Misunderstanding: students conflate "provably equivalent" with "the same
%    program" / "same efficiency", and think ~= says anything about cost.
% -----------------------------------------------------------------------------
\part[4]\hfill

  A student claims: ``if \code{e1} $\eeq$ \code{e2}, then \code{e1} and
  \code{e2} must take the same amount of time to evaluate.''

  Consider these two functions, which you may assume are extensionally
  equivalent (they return the same result on every input):

  \begin{codeblock}
    fun sumTo1 n = if n = 0 then 0 else n + sumTo1 (n - 1)

    fun sumTo2 n = (n * (n + 1)) div 2
  \end{codeblock}

  Is the student's claim true or false? Justify your answer with reference to
  these two functions.

  \begin{solutionorbox}[12em]
    False.

    \code{sumTo1} and \code{sumTo2} are extensionally equivalent --- for every
    \code{n} $\geq 0$ they evaluate to the same integer --- so
    \code{sumTo1 n} $\eeq$ \code{sumTo2 n}. But \code{sumTo1} does $O(n)$
    additions while \code{sumTo2} does a constant number of operations.

    Equivalence ($\eeq$) is a statement about \textbf{what value} an expression
    evaluates to, not \textbf{how much work} it takes to get there. Two
    equivalent expressions may have wildly different work and span.
  \end{solutionorbox}


% -----------------------------------------------------------------------------
% 4. Mutation breaks referential transparency.
%    Misunderstanding: students believe you can always substitute equals for
%    equals (replace `f x` by its "value"), even in the presence of effects.
% -----------------------------------------------------------------------------
\part[4]\hfill

  Consider the following declarations:

  \begin{codeblock}
    val c = ref 0
    fun next () = (c := !c + 1; !c)
  \end{codeblock}

  A student reasons: ``let \code{v = next ()}. Then \code{next () + next ()}
  $\eeq$ \code{v + v} $\eeq$ \code{2 * v}, by substituting the value of
  \code{next ()}.''

  Is this reasoning valid? State the actual value of
  \code{next () + next ()} when it is the first thing evaluated after the
  declarations above, and explain in one sentence what went wrong.

  \begin{solutionorbox}[12em]
    The reasoning is invalid. The actual value is \code{1 + 2} $\eeq$ \code{3},
    not \code{2 * v} for any single \code{v}.

    Each call to \code{next} \textbf{mutates} \code{c} and returns a different
    result, so \code{next ()} is not a fixed value that can be substituted for
    all its occurrences. Mutation breaks \textbf{referential transparency} ---
    the property that lets you replace an expression by its value without
    changing the program's meaning. That substitution is only valid for
    \textit{effect-free} (pure) expressions.
  \end{solutionorbox}
